%% This is file `sample-sigconf-authordraft.tex',
%% generated with the docstrip utility.
%%
%% The original source files were:
%%
%% samples.dtx  (with options: `all,proceedings,bibtex,authordraft')
%% 
%% IMPORTANT NOTICE:
%% 
%% For the copyright see the source file.
%% 
%% Any modified versions of this file must be renamed
%% with new filenames distinct from sample-sigconf-authordraft.tex.
%% 
%% For distribution of the original source see the terms
%% for copying and modification in the file samples.dtx.
%% 
%% This generated file may be distributed as long as the
%% original source files, as listed above, are part of the
%% same distribution. (The sources need not necessarily be
%% in the same archive or directory.)
%%
%%
%% Commands for TeXCount
%TC:macro \cite [option:text,text]
%TC:macro \citep [option:text,text]
%TC:macro \citet [option:text,text]
%TC:envir table 0 1
%TC:envir table* 0 1
%TC:envir tabular [ignore] word
%TC:envir displaymath 0 word
%TC:envir math 0 word
%TC:envir comment 0 0
%%
%%
%% The first command in your LaTeX source must be the \documentclass
%% command.
%%
%% For submission and review of your manuscript please change the
%% command to \documentclass[manuscript, screen, review]{acmart}.
%%
%% When submitting camera ready or to TAPS, please change the command
%% to \documentclass[sigconf]{acmart} or whichever template is required
%% for your publication.
%%
%%
\documentclass[sigconf,authorversion]{acmart}
\usepackage{algorithm}
\usepackage{algorithmic}
\usepackage{xspace}
\usepackage{multirow}

%% comments

\newcommand{\aris}[1]{\noindent\textcolor{orange}{\textbf{[Aris:]}~{#1}}\xspace}
\newcommand{\haoyun}[1]{\noindent\textcolor{blue}{\textbf{[Haoyun:]}~{#1}}\xspace}
\newcommand{\todo}[1]{\noindent\textcolor{magenta}{\textbf{[To do:]}~{#1}}\xspace}

%% paragraphs

\newcommand{\para}[1]{{\noindent{\textbf{#1}}}}
\newcommand{\spara}[1]{{\smallskip\noindent{\textbf{#1}}}}
\newcommand{\mpara}[1]{{\medskip\noindent{\textbf{#1}}}}

%% notation

\newcommand{\reals}{\ensuremath{\mathbb{R}}\xspace}
\newcommand{\bigO}{\ensuremath{\mathcal{O}}\xspace}
\newcommand{\matr}[1]{\ensuremath{\mathbf{{#1}}}\xspace}
\newcommand{\vect}[1]{\ensuremath{\mathbf{{#1}}}\xspace}
\newcommand{\transmatrix}{\ensuremath{\matr{P}}\xspace}
\newcommand{\newtransmatrix}{\ensuremath{\matr{R}}\xspace}
\newcommand{\newtransmatrixopt}{\ensuremath{\newtransmatrix^{*}}\xspace}
\newcommand{\vectx}{\ensuremath{\vect{x}}\xspace}
\newcommand{\vecty}{\ensuremath{\vect{y}}\xspace}
\newcommand{\vectv}{\ensuremath{\vect{v}}\xspace}
\newcommand{\vecte}{\ensuremath{\vect{e}}\xspace}
\newcommand{\vects}{\ensuremath{\vect{s}}\xspace}
\newcommand{\prvector}{\ensuremath{\vect{p}}\xspace}
\newcommand{\indvect}[1]{\ensuremath{\vect{1}_{#1}}\xspace} % indicator vector
\newcommand{\graph}{\ensuremath{G}\xspace}
\newcommand{\vertices}{\ensuremath{V}\xspace}
\newcommand{\edges}{\ensuremath{E}\xspace}
\newcommand{\novertices}{\ensuremath{n}\xspace}
\newcommand{\noedges}{\ensuremath{m}\xspace}
\newcommand{\vertexu}{\ensuremath{u}\xspace}
\newcommand{\vertexv}{\ensuremath{v}\xspace}
\newcommand{\unitary}{\ensuremath{\matr{I}}\xspace}
\newcommand{\matrixU}{\ensuremath{\matr{U}}\xspace}
\newcommand{\matrixE}{\ensuremath{\matr{E}}\xspace}
\newcommand{\nogroups}{\ensuremath{g}\xspace}
\newcommand{\agroup}{\ensuremath{k}\xspace}
\newcommand{\grouppr}{\ensuremath{\phi}\xspace}
\newcommand{\loss}{\ensuremath{L}\xspace}
\newcommand{\lossprev}{\ensuremath{\loss_{\mathit{prev}}}\xspace}
\newcommand{\losscross}{\ensuremath{\loss_{\mathit{cross}}}\xspace}
\newcommand{\feasiblematrices}{\ensuremath{\mathcal{C}}\xspace}
\newcommand{\resfeasiblematrices}{\ensuremath{\mathcal{C}^{\prime}}\xspace}
\newcommand{\boundfactor}{\ensuremath{\delta}\xspace}
\newcommand{\iter}{\ensuremath{{\mathit{iter}}}\xspace}

%% problems

\newcommand{\fpr}{\ensuremath{\phi}{\sc\small PR}\xspace}
\newcommand{\cgfpr}{{\sc\small Cross}-{\fpr}\xspace}
\newcommand{\resfpr}{{\sc\small R}-{\fpr}\xspace}
\newcommand{\rescgfpr}{{\sc\small R}-{\cgfpr}\xspace}

%% algorithms

\newcommand{\fairgd}{{\texttt{FairGD}}\xspace}
\newcommand{\yapprox}{{\texttt{yApprox}}\xspace}
\newcommand{\project}{{\texttt{Project}}\xspace}
\newcommand{\crossgd}{{\texttt{CrossGD}}\xspace}
\newcommand{\pralgo}{{\texttt{PageRank}}\xspace}




%%
%% \BibTeX command to typeset BibTeX logo in the docs
\AtBeginDocument{%
  \providecommand\BibTeX{{%
    Bib\TeX}}}

%% Rights management information.  This information is sent to you
%% when you complete the rights form.  These commands have SAMPLE
%% values in them; it is your responsibility as an author to replace
%% the commands and values with those provided to you when you
%% complete the rights form.
\setcopyright{acmlicensed}
\copyrightyear{2018}
\acmYear{2018}
\acmDOI{XXXXXXX.XXXXXXX}

%% These commands are for a PROCEEDINGS abstract or paper.
\acmConference[Conference acronym 'XX]{Make sure to enter the correct
  conference title from your rights confirmation emai}{June 03--05,
  2018}{Woodstock, NY}
%%
%%  Uncomment \acmBooktitle if the title of the proceedings is different
%%  from ``Proceedings of ...''!
%%
%%\acmBooktitle{Woodstock '18: ACM Symposium on Neural Gaze Detection,
%%  June 03--05, 2018, Woodstock, NY}
\acmISBN{978-1-4503-XXXX-X/18/06}


%%
%% Submission ID.
%% Use this when submitting an article to a sponsored event. You'll
%% receive a unique submission ID from the organizers
%% of the event, and this ID should be used as the parameter to this command.
%%\acmSubmissionID{123-A56-BU3}

%%
%% For managing citations, it is recommended to use bibliography
%% files in BibTeX format.
%%
%% You can then either use BibTeX with the ACM-Reference-Format style,
%% or BibLaTeX with the acmnumeric or acmauthoryear sytles, that include
%% support for advanced citation of software artefact from the
%% biblatex-software package, also separately available on CTAN.
%%
%% Look at the sample-*-biblatex.tex files for templates showcasing
%% the biblatex styles.
%%

%%
%% The majority of ACM publications use numbered citations and
%% references.  The command \citestyle{authoryear} switches to the
%% "author year" style.
%%
%% If you are preparing content for an event
%% sponsored by ACM SIGGRAPH, you must use the "author year" style of
%% citations and references.
%% Uncommenting
%% the next command will enable that style.
%%\citestyle{acmauthoryear}

\newtheorem{problem}{Problem}

%%
%% end of the preamble, start of the body of the document source.
\begin{document}

%%
%% The "title" command has an optional parameter,
%% allowing the author to define a "short title" to be used in page headers.
\title{Fairness-aware PageRank via edge reweighting}

%%
%% The "author" command and its associated commands are used to define
%% the authors and their affiliations.
%% Of note is the shared affiliation of the first two authors, and the
%% "authornote" and "authornotemark" commands
%% used to denote shared contribution to the research.
% \author{Ben Trovato}
% \authornote{Both authors contributed equally to this research.}
% \email{trovato@corporation.com}
% \orcid{1234-5678-9012}
% \author{G.K.M. Tobin}
% \authornotemark[1]
% \email{webmaster@marysville-ohio.com}
% \affiliation{%
%   \institution{Institute for Clarity in Documentation}
%   \streetaddress{P.O. Box 1212}
%   \city{Dublin}
%   \state{Ohio}
%   \country{USA}
%   \postcode{43017-6221}
% }

% \author{Lars Th{\o}rv{\"a}ld}
% \affiliation{%
%   \institution{The Th{\o}rv{\"a}ld Group}
%   \streetaddress{1 Th{\o}rv{\"a}ld Circle}
%   \city{Hekla}
%   \country{Iceland}}
% \email{larst@affiliation.org}

\author{Anonymous Author(s)}
% \email{email1}
% \affiliation{%
%   \institution{institution1}
%   \city{city1}
%  \country{country1}
% }

% \author{author2}
% \email{email2}
% \affiliation{%
%   \institution{institution2}
%   \city{city2}
%   \country{country2}
% }

%%
%% By default, the full list of authors will be used in the page
%% headers. Often, this list is too long, and will overlap
%% other information printed in the page headers. This command allows
%% the author to define a more concise list
%% of authors' names for this purpose.
\renewcommand{\shortauthors}{Trovato et al.}

%%
%% The abstract is a short summary of the work to be presented in the
%% article.
\begin{abstract}
Link-analysis algorithms, such as Page\-Rank, 
are instrumental in understanding the structural dynamics of networks 
by evaluating the importance and influence of individual nodes 
based on their connectivity patterns.
Recently, with the raising importance of responsible AI, 
the question of fairness in link-analysis algorithms has gained traction.

In this paper we present a new approach for incorporating group fairness into the Page\-Rank algorithm,
by reweighting the transition probabilities in the underlying transition matrix.
We formulate the problem of achieving fair Page\-Rank
by seeking to minimize the difference between the actual group-wise Page\-Rank distribution and 
a target fair Page\-Rank distribution. 
% Our problem can be defined for multiple groups. 
We further define a cross-group fairness notion, 
which accounts for group homophily by considering  
random walks with group-biased restart for each group. 
We propose an efficient projected gradient-descent method for computing locally-optimal edge reweightings.
As opposed to earlier approaches, 
we do not recommend new edges to be added in the network, 
nor do we adjust the restart vector, 
instead we keep the topology of the underlying network unchanged
and only modify the relative importance of existing edges. 
We empirically compare our approach with state-of-the-art baselines
and demonstrate the efficacy of our method, 
where very small changes in the transition matrix lead to significant improvement 
in the fairness of the Page\-Rank algorithm.
%
% The new edge weights, computed by our method,  can be used to suggest concrete changes in the network ecosystem, e.g., reweighting and prioritizing the feed of users in a social network.
\end{abstract}

%%
%% The code below is generated by the tool at http://dl.acm.org/ccs.cfm.
%% Please copy and paste the code instead of the example below.
%%
\iffalse
\begin{CCSXML}
<ccs2012>
 <concept>
  <concept_id>00000000.0000000.0000000</concept_id>
  <concept_desc>Do Not Use This Code, Generate the Correct Terms for Your Paper</concept_desc>
  <concept_significance>500</concept_significance>
 </concept>
 <concept>
  <concept_id>00000000.00000000.00000000</concept_id>
  <concept_desc>Do Not Use This Code, Generate the Correct Terms for Your Paper</concept_desc>
  <concept_significance>300</concept_significance>
 </concept>
 <concept>
  <concept_id>00000000.00000000.00000000</concept_id>
  <concept_desc>Do Not Use This Code, Generate the Correct Terms for Your Paper</concept_desc>
  <concept_significance>100</concept_significance>
 </concept>
 <concept>
  <concept_id>00000000.00000000.00000000</concept_id>
  <concept_desc>Do Not Use This Code, Generate the Correct Terms for Your Paper</concept_desc>
  <concept_significance>100</concept_significance>
 </concept>
</ccs2012>
\end{CCSXML}

\ccsdesc[500]{Do Not Use This Code~Generate the Correct Terms for Your Paper}
\ccsdesc[300]{Do Not Use This Code~Generate the Correct Terms for Your Paper}
\ccsdesc{Do Not Use This Code~Generate the Correct Terms for Your Paper}
\ccsdesc[100]{Do Not Use This Code~Generate the Correct Terms for Your Paper}
\fi 

%%
%% Keywords. The author(s) should pick words that accurately describe
%% the work being presented. Separate the keywords with commas.
\keywords{Link analysis, Page\-Rank, graph analysis, fairness}
%% A "teaser" image appears between the author and affiliation
%% information and the body of the document, and typically spans the
%% page.
% \begin{teaserfigure}
%   \includegraphics[width=\textwidth]{sampleteaser}
%   \caption{Seattle Mariners at Spring Training, 2010.}
%   \Description{Enjoying the baseball game from the third-base
%   seats. Ichiro Suzuki preparing to bat.}
%   \label{fig:teaser}
% \end{teaserfigure}

\received{20 February 2007}
\received[revised]{12 March 2009}
\received[accepted]{5 June 2009}

%%
%% This command processes the author and affiliation and title
%% information and builds the first part of the formatted document.
\maketitle

\iffalse
\aris{I suggest using macros for all the math variables in the paper. 
In this way we can change the notation easily, if needed. 
For example, see the \texttt{macros.tex} file, 
where I have defined some macros and have used them in the main file.}
\fi

\iffalse
\aris{Do not use citations as proper nouns. 
For instance, it is OK to write
``The Page\-Rank algorithm \cite{brin_anatomy_1998} ...''
but not OK to write 
``... adapted from \cite{tsioutsiouliklis_fairness-aware_2021}.''
Instead we should write 
``... adapted from \citet{tsioutsiouliklis_fairness-aware_2021}.''
A good idea is to use only commands
\texttt{$\backslash$citet} (when citation is used as noun) and
\texttt{$\backslash$citep} (when citation is not used as noun), 
and never use \texttt{$\backslash$cite}.}
\fi 

\section{Introduction [Aris]}

Machine-learning algorithms and data-driven decision making are playing 
increasingly crucial roles in our lives, 
impacting job markets, healthcare, government, industry, and education. 
However, as reliance on automated decisions grows, 
so does the concern for fairness in algorithms. 
Ensuring that machine-learning algorithms are reliable
and satisfy strong notions of fairness
has become an issue of societal concern and, 
justifiably, has received significant attention in the computer-science literature.
Algorithmic fairness has been studied for
many different tasks of primarily supervised~\citep{hardt2016equality,joseph2016fairness,pedreshi2008discrimination,zemel2013learning},
but also unsupervised learning~\citep{bera2019fair,chen2019proportionally,chierichetti2017fair}.
%
In the context of network algorithms, 
which is the focus of this paper, 
a few different approaches for fairness-aware link analysis 
have been proposed recently~\citep{khajehnejad_crosswalk_2022,rahman_fairwalk_2019,tsioutsiouliklis_fairness-aware_2021,tsioutsiouliklis_link_2022,zeng_fair_2021}, 
however, the topic has received significantly less attention
than the study of fairness in other machine-learning problems. 

Link-analysis algorithms, such as PageRank, 
are fundamental tools for analyzing large-scale network data. 
By quantifying the importance or relevance of nodes within a network, 
based on their connections and the quality of those connections, 
these algorithms provide a means to identify authoritative nodes, detect communities, 
uncover patterns of influence, and mitigate vulnerabilities within networks. 
Link-analysis algorithms have been used in numerous applications, 
including, web search, social-network analysis, and recommendation systems.

In this paper we present a new approach for incorporating fairness into the Page\-Rank algorithm,
by reweighting the transition probabilities of the Page\-Rank matrix. 
More concretely, we formulate the problem of minimizing the Page\-Rank distance
to a fair distribution, while allowing only for reweighting the network edges
within prespecified limits. 
With our formulation we can quantify and improve fairness 
with respect to many vertex groups in the network.
As opposed to earlier approaches, 
we do not recommend new edges to be added in the network, 
nor do we adjust the restart vector, 
instead we keep the topology of the underlying network unchanged
and only modify the relative importance of existing edges. 
If perfect fairness cannot be achieved with the permissible transition-matrix modifications, 
our formulation aims to find a solution that is as close as possible 
to the target fairness objective.
Our approach provides an actionable environment for achieving fairness in Page\-Rank-type methods, 
as the new edge weights, computed by our method, 
can be used to suggest concrete changes in the network ecosystem,
e.g., reweighting and prioritizing the feed of users in a social network. 

We propose a projected gradient descent method for computing locally-optimal edge reweightings
and we show how it can be implemented efficiently. 
We empirically compare our approach with state-of-the-art baselines
and demonstrate the efficacy of our method, 
where very small changes in the transition matrix lead to significant improvement 
in the fairness of the Page\-Rank algorithm.


\citet{tsioutsiouliklis_fairness-aware_2021} measure the importance of a group by the total Page\-Rank scores of vertices in the group. They rebalance the graph such that the Page\-Rank vector is $\grouppr$-fair, i.e., the total Page\-Rank scores of vertices in the protected group is $\grouppr$. We extend this work to multiple groups. Additionally, instead of having the total Page\-Rank scores in the groups exactly equal to $\grouppr$, we would like to approximate the total Page\-Rank scores in the groups to $\grouppr$. 

\aris{List of contributions below: 
(Haoyun: feel free to edits the contributions, 
and list additional contributions, as you consider appropriate.)
\begin{itemize}
\item New fairness definition of Page\-Rank, based on matrix reweighting and 
        not changing the network topology.
\item minimize the loss function with respect to target fairness
\item notion extends to $k$ groups, and cross-group fairness
\item gradient descent method converges to local optimum
\item intuitive interpretation of gradient
\item efficient computation via approximation of power iteration
\item experimentally, method outperfroms baselines and makes very small changes in the transition matrix
\end{itemize}
}

\section{Related Work [Haoyun]}

\aris{haoyun: please add some related papers + references and one-sentence description of each paper.}

\textbf{FairWalk}~\citep{rahman_fairwalk_2019}. In node embedding learning, when choosing the next step in a random walk, equal probability is assigned to all groups in the neighborhood of the current node.

\textbf{CrossWalk}~\citep{khajehnejad_crosswalk_2022}. In addition to FairWalk, CrossWalk further reweights edges according to the distance to the group boundary.

\textbf{Fairness-aware Page\-Rank}. \citet{tsioutsiouliklis_fairness-aware_2021} defines a novel metric for fairness in Page\-Rank and proposes two families of algorithms FSPR and LFPR to achieve this fairness. FSPR changes the restart vector, and LFPR reweights the transition matrix with introducing new edges.

\textbf{Link Recommendations for Page\-Rank Fairness}. \citet{tsioutsiouliklis_link_2022} analyzes the effect of potentially new edges as well as existed edges on the fairness of Page\-Rank. They proposes a linear time link recommendation algorithm that maximizes fairness.

\textbf{Fair representation learning for heterogeneous information networks}. \citet{zeng_fair_2021} chooses vertex in the disadvantaged groups with higher probability when generating meta-paths.

\textbf{Applying Fairness Constraints on Graph Node Ranks Under Personalization Bias}. \citet{krasanakis_applying_2021} mitigates disparate impact in ranking under personalization bias by a personalization editing mechanism.

\section{Preliminaries and problem definition}

% \aris{Let consistency let us use either ``vertex'' or ``node'', not both. I suggest ``vertex''.}
% \haoyun{Changed ``node`` to ``vertex``.}

We consider a directed graph $\graph = (\vertices, \edges)$,  
where $\vertices$ is a set of vertices, and $\edges$ is a set of edges. 
We denote by $\novertices = |\vertices|$ the number of vertices and 
by $\noedges = |\edges|$ the number of edges. 
The Page\-Rank algorithm assigns a score to each vertex in the graph~\citep{brin_anatomy_1998}. 
The Page\-Rank score of a vertex \vertexu equals to the probability 
in the stationary distribution of a random walk on the graph \graph 
visiting the vertex~\vertexu.

In the following we consider that the graph \graph is 
represented by a row-stochastic transition matrix \transmatrix with nonnegative entries.
The transition matrix \transmatrix is derived by normalizing each row of the adjacency matrix of the graph, 
and replacing zero rows (corresponding to vertices with out-degree equal to 0) 
with the restart vector \vectv, which is introduced below.
Page\-Rank is a widely-studied algorithm and there are many equivalent formulations. 
Here we adopt the notation from the work of \citet{tsioutsiouliklis_fairness-aware_2021}, 
who introduced the idea of fairness-aware Page\-Rank.

% \aris{I suggest denoting the Page\-Rank vector using a different letter than \vectx,  
% for example \prvector.  This is because \vectx is most often used as a variable name.}

\begin{definition}
Given a nonnegative row-stochastic transition matrix \transmatrix, 
a restart probability $0 < \gamma < 1$, and a restart vector \vectv with $||\vectv||_1 = 1$, 
the Page\-Rank score vector \prvector is the solution to the following equation:
\begin{equation}
\label{eq:pagerank}
    \prvector^T = (1 - \gamma) \prvector^T \transmatrix + \gamma \vectv^T.
\end{equation}
\end{definition}

Solving Equation~\eqref{eq:pagerank} gives
\begin{equation}
    \prvector^T = \gamma \vectv^T (\unitary - (1 - \gamma) \transmatrix)^{-1}.
\end{equation}

We hence view the Page\-Rank score vector \prvector as a function of the transition matrix \transmatrix, 
the restart probability $\gamma$, and the restart vector~\vectv.

\subsection{Fair Page\-Rank}

Similarly to the earlier definition of fairness-aware Page\-Rank~\citep{tsioutsiouliklis_fairness-aware_2021},
we consider the set of vertices in the network being partitioned into distinct groups. 
Those can represent, for instance, social groups defined using 
sensitive attributes, e.g., gender, ethnicity, sexual orientation, etc.
While earlier studies on fairness-aware Page\-Rank have considered only two groups, 
our method can easily extend to a larger number of groups. 
In particular, we assume that the vertex set $\vertices$ is partitioned 
into \nogroups groups $\vertices_1, \ldots, \vertices_\nogroups$, 
for some $\nogroups \ge 2$.

To operationalize the notion of Page\-Rank fairness we assume 
a \emph{target Page\-Rank score} $\grouppr_\agroup$,
for each vertex group $\agroup=1,\ldots,\nogroups$, 
which is the total Page\-Rank mass expected to be assigned by the Page\-Rank algorithm
on the vertices of group $\vertices_\agroup$.
%
A natural choice is to set $\grouppr_\agroup$
equal to the fraction of vertices in~$\vertices_\agroup$ within~\vertices, 
i.e., $\grouppr_\agroup = |\vertices_\agroup| / \novertices$.
However, other choices are possible depending on the application scenario. 
In this work, as in earlier papers, 
we assume that the target Page\-Rank score $\grouppr_\agroup$, 
for each group $\agroup=1,\ldots,\nogroups$,
is provided as input to our problem setting.

Let us denote the indicator vector of group $\vertices_\agroup$ as 
$\indvect{\agroup} \in \{0, 1\}^{\novertices}$,
that is, $\indvect{\agroup}[i] = 1$ if vertex $i \in \vertices_\agroup$ and $\indvect{\agroup}[i] = 0$ otherwise. 
Consider the Page\-Rank vector \prvector
computed by the Page\-Rank algorithm on a transition matrix \transmatrix, 
with restart probability~$\gamma$, and restart vector~\vectv.
The total Page\-Rank score in group $\vertices_\agroup$ is given by $\indvect{\agroup}^T \prvector$. 

Given an input graph \graph, and associated transition matrix \transmatrix, 
our objective is to find a new transition matrix $\newtransmatrix$, 
defined over the same input graph \graph, 
so that the group-wise PageRank distribution computed over \newtransmatrix (to be defined shortly)
is as close as possible to the group-wise target 
PageRank distribution $\{\grouppr_\agroup\}_{\agroup=1,\ldots,\nogroups}$.

Implicit in our formulation is the idea that it is meaningful to 
modify the edge weights of a graph 
in order to nudge the PageRank scores into a more fair distribution. 

As a simple example, given two entries of the transition matrix $\transmatrix[u,v]=\transmatrix[u,w]=0.5$,
indicating that a random walk from vertex~$u$ proceeds to vertices $v$ or $w$ with probability $0.5$ each, 
a permissible solution could be a new transition matrix \newtransmatrix 
with entries $\newtransmatrix[u,v]=0.65$ and $\newtransmatrix[u,w]=0.35$,
indicating a random walk with transition from $u$ to $v$ or $w$ with probabilities $0.65$ and $0.35$, respectively.

To preserve the original structure of the input graph \graph 
we require that the new transition matrix $\newtransmatrix$ 
satisfies $\newtransmatrix[i, j] = 0$ if $\transmatrix[i, j] = 0$. 
This condition ensures that no new edges are added
to the graph that correspond to the new transition matrix $\newtransmatrix$.

To formalize our problem, 
we first define the fairness loss function.
\begin{definition}
\label{definition:fairness-loss}
Given a row-stochastic transition matrix \transmatrix, 
a restart probability $0 < \gamma < 1$, 
a restart vector $\vectv \in \reals^{\novertices}$ with $\|\vectv\|_1 = 1$, 
groups $\vertices_1, \ldots, \vertices_\nogroups$, and 
group-wise target Page\-Rank scores $\grouppr_1, \ldots, \grouppr_\nogroups$, 
for $\nogroups\ge 2$, 
the \emph{fairness loss function} is defined as
\begin{equation}
    \loss(\transmatrix, \gamma, \vectv) = 
        \frac{1}{\nogroups} \sum_{\agroup=1}^\nogroups \left(\indvect{\agroup}^T \prvector - \grouppr_\agroup\right)^2,
\label{equation:fairness-loss}
\end{equation}
where \prvector is the Page\-Rank vector with respect to parameters $(\transmatrix, \gamma, \vectv)$.
\end{definition}

We consider that a Page\-Rank vector \prvector is closer to $\grouppr$-fairness 
if the fairness loss is smaller. 
Moreover, a Page\-Rank vector \prvector is $\grouppr$-fair 
if it incurs zero fairness loss. 
We can now formulate the $\grouppr$-fair Page\-Rank (\fpr) problem as follows:

\begin{problem}[$\grouppr$-fair Page\-Rank (\fpr)]
Given a graph $\graph = (\vertices, \edges)$ with row-stochastic transition matrix \transmatrix, 
a restart probability $0 < \gamma < 1$, 
a restart vector $\vectv \in \reals^{\novertices}$ with $||\vectv||_1 = 1$, 
groups $\vertices_1, \ldots, \vertices_\nogroups$, and 
group-wise target Page\-Rank scores $\grouppr_1, \ldots, \grouppr_\nogroups$, 
for $\nogroups\ge 2$, 
find a new transition matrix $\newtransmatrixopt$
that minimizes the fairness loss function, i.e.,
\begin{equation}
\newtransmatrixopt = \arg \min_{\newtransmatrix \in \feasiblematrices(\transmatrix)} \loss(\newtransmatrix, \gamma, \vectv),
\end{equation}
where $\feasiblematrices(\transmatrix)$ is the set of row-stochastic matrices 
that have non-negative entries and satisfy 
$\newtransmatrix[i, j] = 0$ if $\transmatrix[i, j] = 0$.
\end{problem}

% \aris{say here that we change only \transmatrix, not $\gamma$ nor \vectv),}

\todo{Show that the loss function is not convex.}

\subsection{Cross-group fair Page\-Rank}

In the \fpr %  $\grouppr$-fair Page\-Rank 
problem we require the total Page\-Rank score $\indvect{\agroup}^T\prvector$ 
in each group $\vertices_\agroup$ to be equal to $\grouppr_\agroup$. 
When it is not possible to satisfy the requirement $\indvect{\agroup}^T \prvector = \grouppr_\agroup$
using only edge re\-weightings, 
our formulation seeks to minimize the fairness loss function $\loss(\cdot, \gamma, \vectv)$,
so that the group-wise Page\-Rank scores $\indvect{\agroup}^T \prvector$
are as close as possible to the corresponding target scores $\grouppr_\agroup$.

For computing the Page\-Rank vector \prvector 
we assume that the restart probability $\gamma$ and the restart vector \vectv are given as input. 
In the absence of a preferred restart vector it is common to consider the uniform vector
$\vectv_0 = \frac{1}{\novertices} [1,\ldots,1]$.
On the other hand, real-world networks are mostly homophilic, 
i.e., vertices in the same group are more likely to be connected to each other. 
Hence, the total Page\-Rank scores in group $\vertices_\agroup$ are largely contributed 
by random walks (re)\-starting at vertices in group $\vertices_\agroup$. 
%
To increase the diversity in the network, 
it is desirable for the total Page\-Rank score $\indvect{\agroup}^T\prvector$  
in group~$\vertices_\agroup$ to be equal to~$\grouppr_\agroup$
\emph{when random walks (re)\-start at any~group}. 

In this section, we propose the notion of \emph{cross-group fairness}, 
which enhances fair Page\-Rank 
by considering random walks that (re)\-start at each group, 
and penalizing the fairness loss for each such (re)\-starting batch.

More specifically, 
let $\vectv_\agroup = \frac{1}{|\vertices_\agroup|} \indvect{\agroup}$ 
be the normalized indicator vector of group $\vertices_\agroup$.
We define the \emph{cross-group fairness loss function} as follows:

\begin{definition}
\label{definition:cgfl}
Given a row-stochastic transition matrix \transmatrix, 
a restart probability $0 < \gamma < 1$, 
groups $\vertices_1, \ldots, \vertices_\nogroups$, and 
group-wise target Page\-Rank scores $\grouppr_1, \ldots, \grouppr_\nogroups$, 
for $\nogroups \ge 2$,
the cross-group fairness loss function is defined as 
\begin{equation}
\label{equation:cgfl}
    \losscross(\transmatrix, \gamma)  
        = \frac{1}{\nogroups} \sum_{\ell=1}^\nogroups \loss(\transmatrix, \gamma, \vectv_{\ell})
        = \frac{1}{\nogroups^2} \sum_{\ell=1}^\nogroups \sum_{\agroup=1}^\nogroups 
                \left(\indvect{\agroup}^T \prvector_\ell - \grouppr_\agroup\right)^2,
\end{equation}
where $\prvector_\ell$ is the Page\-Rank vector with respect to parameters 
$(\transmatrix, \gamma, \vectv_{\ell})$.
\end{definition}

In the cross-group fairness loss, 
we iterate through all groups and sum up the fairness loss
when restarting the random walk at each group. 
Asking to minimize cross-group fairness loss function
suggests the cross-group $\grouppr$-fair Page\-Rank (\cgfpr) problem.

\begin{problem}[Cross-Group $\grouppr$-fair Page\-Rank (\cgfpr)]
Given a graph $\graph = (\vertices, \edges)$ with  
row-stochastic transition matrix \transmatrix, 
a restart probability $0 < \gamma < 1$, 
groups $\vertices_1, \ldots, \vertices_\nogroups$, and 
group-wise target Page\-Rank scores $\grouppr_1,  \ldots, \grouppr_\nogroups$, 
for $\nogroups\ge 2$, 
find a new transition matrix $\newtransmatrixopt$ that minimizes the cross-group fairness loss function, i.e.,
\begin{equation}
    \newtransmatrixopt = \arg\min_{\newtransmatrix \in \feasiblematrices(\transmatrix)} 
                            \losscross(\newtransmatrix, \gamma),
\end{equation}
where $\feasiblematrices(\transmatrix)$ is the set of row-stochastic matrices 
that have non-negative entries and satisfy $\newtransmatrix[i, j] = 0$ if $\transmatrix[i, j] = 0$.
\end{problem}

\subsection{Restricting edge reweighting}

In the \fpr and \cgfpr problems, discussed in the previous section, 
we require that the changes in the transition matrix \transmatrix, 
to obtain a new transition matrix \newtransmatrix, 
should take place only on existing edges.
However, we do not restrict the \emph{amount} of modification on those edges.
In some cases, the weight of an existing edge could be reduced to zero,
which is equivalent to removing the edge from the graph.
To better preserve the original structure of the input graph \graph,
we restrict the element-wise changes in the transition matrix \transmatrix
by introducing lower- and upper-bound constraints on 
the entries of the new transition matrix \newtransmatrix.
Applying this restriction to the \fpr problem,
we have the following restricted $\grouppr$-fair Page\-Rank (\resfpr)~problem.

\begin{problem}[Restricted $\grouppr$-fair Page\-Rank (\resfpr)]
Given a graph $\graph = (\vertices, \edges)$ with row-stochastic transition matrix \transmatrix, 
a restart probability $0 < \gamma < 1$, 
a restart vector $\vectv \in \reals^{\novertices}$ with $||\vectv||_1 = 1$, 
groups $\vertices_1, \ldots, \vertices_\nogroups$, 
group-wise target Page\-Rank scores $\grouppr_1, \ldots, \grouppr_\nogroups$, 
for $\nogroups\ge 2$, 
and a modification-bound factor $\boundfactor \ge 0$,
find a new transition matrix $\newtransmatrixopt$
that minimizes the fairness loss function, i.e.,
\begin{equation*}
  \newtransmatrixopt = \arg \min_{\newtransmatrix \in \resfeasiblematrices(\transmatrix)} \loss(\newtransmatrix, \gamma, \vectv),
\end{equation*}
where $\resfeasiblematrices(\transmatrix)$ is the set of row-stochastic matrices
that have non-negative entries, and for all $i, j$ satisfy the constraint
% 
% $\newtransmatrix[i, j] = 0$ if $\transmatrix[i, j] = 0$,
% and additionally
\begin{equation}
\label{inequalities:change-constraint}
    (1 - \boundfactor) \transmatrix[i, j] \le \newtransmatrix[i, j] \le (1 + \boundfactor) \transmatrix[i, j] .
\end{equation} 
\end{problem}

In other words, Inequalities~(\ref{inequalities:change-constraint})
restrict the entries of the new transition matrix \newtransmatrix
to be within factors $(1 - \boundfactor)$ and $(1 + \boundfactor)$
of the entries of the original transition matrix \transmatrix.
Notice that Inequalities~(\ref{inequalities:change-constraint})
imply also the constraint ``if $\transmatrix[i, j] = 0$ then $\newtransmatrix[i, j] = 0$,''
so it is not necessary to specify that constraint separately. 

Similarly, we define the restricted cross-group $\grouppr$-fair Page\-Rank (\rescgfpr) problem.

\begin{problem}[Restricted Cross-Group $\grouppr$-fair Page\-Rank \\ (\rescgfpr)]
Given a graph $\graph = (\vertices, \edges)$ with  
row-stochastic transition matrix \transmatrix, 
a restart probability $0 < \gamma < 1$, 
groups $\vertices_1, \ldots, \vertices_\nogroups$, 
group-wise target Page\-Rank scores $\grouppr_1,  \ldots, \grouppr_\nogroups$, 
for $\nogroups\ge 2$, 
and a modification-bound factor $\boundfactor \ge 0$,
find a new transition matrix $\newtransmatrixopt$ that minimizes the cross-group fairness loss function, i.e.,
\begin{equation*}
    \newtransmatrixopt = \arg\min_{\newtransmatrix \in \resfeasiblematrices(\transmatrix)} 
                         \losscross(\newtransmatrix, \gamma),
\end{equation*}
where $\resfeasiblematrices(\transmatrix)$ is the set of row-stochastic matrices
that have non-negative entries, and for all $i, j$ satisfy the constraint
\begin{equation*}
    (1 - \boundfactor) \transmatrix[i, j] \le \newtransmatrix[i, j] \le (1 + \boundfactor) \transmatrix[i, j] .
\end{equation*} 
\end{problem}

\section{Algorithms}

In this section, we propose a projected gradient-descent algorithm 
for finding the optimal transition matrix $\newtransmatrix$ for the $\grouppr$-fair Page\-Rank problem. 
We then extend the algorithm to solve the cross-group $\grouppr$-fair Page\-Rank problem.

\subsection{Projected gradient descent for \grouppr-fair Page\-Rank}

We first derive the gradient of the fairness loss function 
$\loss(\transmatrix, \gamma, \vectv)$ 
with respect to the transition matrix \transmatrix.
\begin{theorem}
\label{theorem:gradient-fairness-loss}
The gradient of the fairness loss function 
$\loss(\transmatrix, \gamma, \vectv)$ 
with respect to the transition matrix \transmatrix is given by
\begin{equation*}
% \label{eq:thm_gradient_fairness_loss}
\begin{aligned}
    \frac{\partial}{\partial \transmatrix} \loss(\transmatrix, \gamma, \vectv) 
        & = 2 \gamma (1 - \gamma) \sum_{k=1}^{\nogroups} 
            \left(\indvect{\agroup}^T \prvector - \grouppr_\agroup\right) 
            \matrixU^{-1} \vectv \indvect{\agroup}^T \matrixU^{-1} \\
        & = \frac{2(1 - \gamma)}{\gamma} \sum_{k=1}^{\nogroups} 
            \left(\indvect{\agroup}^T \prvector - \grouppr_\agroup\right) \prvector \vecty_\agroup^T,
\end{aligned}
\end{equation*}
where $\matrixU = \unitary - (1 - \gamma) \transmatrix^T$, 
\prvector is the Page\-Rank vector with respect to $(\transmatrix, \gamma, \vectv)$, and 
$\vecty_\agroup = \gamma \matrixU^{-T} \indvect{\agroup} = 
    \gamma (\unitary - (1 - \gamma) \transmatrix)^{-1} \indvect{\agroup}$.
\end{theorem}
  
\begin{proof}
Recalling that $\prvector = \gamma \matrixU^{-1} \vectv$,
we have the following:
\begin{equation}
\begin{aligned}
\label{eq:gradient_fairness_loss}
    \frac{\partial}{\partial \transmatrix} \loss(\transmatrix, \gamma, \vectv) 
        & = \frac{\partial}{\partial \transmatrix} 
                \sum_{k=1}^{\nogroups} \left(\indvect{\agroup}^T \prvector - \grouppr_\agroup\right)^2 \\
        & = 2 \sum_{k=1}^{\nogroups} \left(\indvect{\agroup}^T \prvector - \grouppr_\agroup\right) 
                \frac{\partial}{\partial \transmatrix} \left(\indvect{\agroup}^T \prvector - \grouppr_\agroup\right) \\
        & = 2 \sum_{k=1}^{\nogroups} \left(\indvect{\agroup}^T \prvector - \grouppr_\agroup\right) 
                \frac{\partial}{\partial \transmatrix} \left(\indvect{\agroup}^T \prvector\right)  \\ 
        & = 2 \gamma \sum_{k=1}^{\nogroups} \left(\indvect{\agroup}^T \prvector - \grouppr_\agroup\right) 
                \frac{\partial}{\partial \transmatrix} \left(\indvect{\agroup}^T \matrixU^{-1} \vectv\right). 
\end{aligned}
\end{equation}
By the chain rule of matrix calculus \citep[Equation~(126)]{petersen_matrix_2008}, 
we have
\begin{equation}
\label{eq:chain_rule}
\begin{aligned}
\frac{\partial}{\partial \transmatrix[i, j]} \left(\indvect{\agroup}^T \matrixU^{-1} \vectv\right) 
    & = \mathrm{Tr} 
        \left[ 
            \left( \frac{\partial (\indvect{\agroup}^T \matrixU^{-1} \vectv)}{\partial \matrixU} \right)^T 
            \frac{\partial \matrixU}{\partial \transmatrix[i, j]} 
        \right]
\end{aligned}
\end{equation}
By considering the derivative of the inverse~\cite[Equation~(55)]{petersen_matrix_2008}, it is

\begin{equation*}
    \frac{\partial (\indvect{\agroup}^T \matrixU^{-1} \vectv)}{\partial \matrixU} 
        = -\matrixU^{-T} \indvect{\agroup} \vectv^T \matrixU^{-T}.
\end{equation*}
We also have 
\begin{equation*}
    \frac{\partial \matrixU}{\partial \transmatrix[i, j]} 
        = - (1 - \gamma) \frac{\partial \transmatrix^T}{\partial \transmatrix[i, j]} = - (1 - \gamma) \matrixE_{ij},
\end{equation*}
where $\matrixE_{ij}$ is the matrix whose $[i,j]$ entry is 1 and all other entries are 0.
From Equation~\eqref{eq:chain_rule}, we have
\begin{equation}
\label{eq:gradient_onek_Uinv_v}
\begin{aligned}
\frac{\partial}{\partial \transmatrix[i, j]} \left(\indvect{\agroup}^T \matrixU^{-1} \vectv\right) 
    & = (1 - \gamma) \mathrm{Tr} 
            \left[ \left( \matrixU^{-T} \indvect{\agroup} \vectv^T \matrixU^{-T} \right)^T \matrixE_{ij} \right] \\ 
    & = (1 - \gamma) \left( \matrixU^{-1} \vectv \indvect{\agroup}^T \matrixU^{-1} \right)[i, j] \\
    & = \frac{1 - \gamma}{\gamma^2} \left( \prvector \vecty_\agroup^T \right)[i, j].
\end{aligned}
\end{equation}
Substituting~\eqref{eq:gradient_onek_Uinv_v} into Equation~\eqref{eq:gradient_fairness_loss} 
completes the proof.
\end{proof}

In Theorem~\ref{theorem:gradient-fairness-loss} we have defined
$\vecty_\agroup = \gamma \matrixU^{-T} \indvect{\agroup}$.
The $j$-th coordinate of~$\vecty_\agroup$ is equal to 
$\vecty_{\agroup}[j] = \indvect{\agroup}^T (\gamma \matrixU^{-1} \mathbf{e}_j)$
and is the sum of Page\-Rank scores of vertices in group $\vertices_\agroup$
with restart at the one-hot vector~$\vecte_j$, 
i.e., the random walk always (re)\-starts at vertex $j$. 

Hence, the gradient of the fairness loss has the following intuitive interpretation:
the $[i, j]$ entry of the gradient matrix is a weighted sum of the products of 
$\prvector[i]$ and $\vecty_\agroup[j]$. 
The coordinate $\prvector[i]$ represents the frequency that a random walk 
restarting at \vectv visits vertex~$i$, 
while the coordinate $\vecty_\agroup[j]$ represents the frequency that a random walk restarting at vertex~$j$ 
visits vertices in the group~$\vertices_\agroup$. 
Hence, the $\agroup$-th term of the summation represents the weighted contribution of the edge $(i, j)$ 
on the total Page\-Rank score of group~$\vertices_\agroup$.
The weight of the $\agroup$-th term is proportional to the difference between the total Page\-Rank scores 
in group $\vertices_\agroup$ and the group-wise target score~$\grouppr_\agroup$.

In terms of computational cost, 
the gradient of the fairness loss 
involves the inverse of the matrix $\matrixU = \unitary - (1 - \gamma) \transmatrix^T$, 
which is computationally expensive. 
To avoid a matrix inversion, 
we approximate $\vecty_\agroup = \gamma \matrixU^{-T} \indvect{\agroup}$ with a power iteration:
\begin{equation}\label{eq:power_iteration}
  \begin{aligned}
    \vecty_\agroup 
        & = \gamma \sum_{i=0}^{\infty} (1 - \gamma)^i \transmatrix^i \indvect{\agroup} \\
        & = \gamma \left(\indvect{\agroup} + (1 - \gamma) \transmatrix \indvect{\agroup} + (1 - \gamma)^2 \transmatrix (\transmatrix \indvect{\agroup}) + \ldots\right).
  \end{aligned}
\end{equation}
Denoting by \newtransmatrix the transition matrix sought by the projected gradient-descent method, 
after each gradient descent step, 
\newtransmatrix is projected back to the feasible set 
$\feasiblematrices(\transmatrix)$ by the following steps:
\begin{enumerate}
  \item For all entries $[i, j]$, if $\transmatrix[i, j] = 0$ then $\newtransmatrix[i, j] = 0$.
  \item All negative entries in $\newtransmatrix$ are set to zero.
  \item All rows of \newtransmatrix are normalized, so that \newtransmatrix is row-stochastic.
\end{enumerate}

The first step of the projection is skipped 
in the inplementation, 
since we only store and update matrix entries where $\transmatrix[i, j] > 0$. 

\haoyun{I am not sure if it is necessary to discuss how to handle sink vertices. \\
1. It is not the emphasis of this paper, \\
2. There is almost no change in results with or without handling the sink vertices. \\
3. I find it difficult to discuss it clearly.}

\aris{I am not sure what you mean by 
\emph{``Hence, we ignore the sink vertices in \transmatrix. This will not affect the Page\-Rank algorithm.''}
Can you elaborate a bit? \\
Also, what happens if there is the gradient-descent creates a new row with zeros
(due to step (2) of the projection)?
}

\haoyun{I meant Page\-Rank by default handles sink vertices by replacing the zero rows 
in the transition matrix with the restart vector, 
so that we can ignore the sink vertices in the transition matrix. \\
Now I rewrite the paragraph and make it clearer. \\
I didn't think about the case where the gradient-descent creates a zero row.
It is too rare to happen, but if it happens, we can replace the zero row with a uniform vector.}


Since rows in the transition matrix \transmatrix corresponding to sink vertices 
are replaced by the restart vector \vectv, every sink vertices will contribute 
$\bigO(\novertices)$ non-zero entries in the transition matrix \transmatrix. 
To avoid a large number of non-zero entries in the transition matrix \transmatrix
because of sink vertices in order to reduce the computational cost,
we decompose the transition matrix \transmatrix into a matrix $\transmatrix_{0}$
considering transition only for non-sink vertices,
and another matrix $\transmatrix_{s}$ considering transition only for sink vertices.
That is,
\begin{equation}
  \transmatrix = \transmatrix_{0} + \transmatrix_{s} = \transmatrix_{0} + \indvect{s} \vectv^T, 
\end{equation}
with $\indvect{s}[i]=1$ if vertex $i$ is a sink vertex and otherwise $indvect{s}[i]=0$.
Therefore we can store and update only $\bigO(m)$ non-zero entries of $\transmatrix_{0}$.

Equation~\eqref{eq:pagerank} is then modified to
\begin{equation}
    \prvector^T = (1 - \gamma) \prvector^T (\transmatrix_{0} + \indvect{s} \vectv^T) + \gamma \vectv^T,
\end{equation}
and Equation~\eqref{eq:power_iteration} is modified to
\begin{equation}
%  \begin{aligned}
    \vecty_\agroup = \gamma \sum_{i=0}^{\infty} (1 - \gamma)^i (\transmatrix_{0} + \indvect{s} \vectv^T)^i \indvect{\agroup}.
%         & = \gamma (\indvect{\agroup} + (1 - \gamma) \transmatrix_{\text{{true}}} \indvect{\agroup} + (1 - \gamma)^2 \transmatrix_{\text{true}} (\transmatrix_{\text{true}} \indvect{\agroup}) + \ldots).
%  \end{aligned}
\end{equation}

To sum up, our method is presented in Algorithms~\ref{alg:fairgd},~\ref{alg:poweriteration}, and~\ref{alg:project}.
Algorithm~\ref{alg:fairgd} (\fairgd) 
is the overall projected gradient-descent scheme that 
iteratively updates the transition matrix \newtransmatrix until convergence.
Algorithm~\ref{alg:poweriteration} (\yapprox)
presents the approximation of vector $\vecty_\agroup$
using power iterations.
Algorithm~\ref{alg:project} (\project) is the projection function, 
where the transition matrix \newtransmatrix is projected back to the feasible set after each iteration.

\begin{algorithm}[t]
  \caption{\label{alg:fairgd}\fairgd: fair Page\-Rank projected gradient descent}
  \begin{algorithmic}[1]
    \REQUIRE Transition matrix~\transmatrix, restart probability~$\gamma$, restart vector~\vectv, 
    groups $\vertices_1, \ldots, \vertices_\nogroups$, target Page\-Rank scores $\grouppr_1, \ldots, \grouppr_\nogroups$, learning rate~$\alpha$, power iteration exponent~$t$, convergence criterion $\epsilon$
    \ENSURE Optimized transition matrix $\newtransmatrix$
    \STATE $\newtransmatrix \gets \transmatrix$
    \STATE ${\iter} \gets 0$ \aris{\iter is not used}
    \STATE ${\lossprev} \gets \infty$
    \WHILE {$|\loss(\newtransmatrix, \gamma, \vectv) - {\lossprev}| > \epsilon$}
      \STATE $\prvector \gets {\pralgo}(\newtransmatrix, \gamma, \vectv)$
      \STATE ${\lossprev} \gets \loss(\newtransmatrix, \gamma, \vectv)$
      \FOR {$\agroup = 1$ to $\nogroups$}
        \STATE $\theta_\agroup \gets \indvect{\agroup}^T \prvector$
        \STATE $\vecty_{\agroup} \gets$ ${\yapprox}(\newtransmatrix, \gamma, \indvect{\agroup}, t)$
        \FOR {$(i, j)$ where $\transmatrix[i, j] > 0$}
          \STATE $\newtransmatrix[i, j] \gets \newtransmatrix[i, j] - \alpha \frac{2(1 - \gamma)}{\gamma} (\theta_\agroup - \grouppr_\agroup) \prvector[i] \vecty_{\agroup}[j]$
        \ENDFOR
        \ENDFOR
        \STATE $\newtransmatrix \gets {\project}(\newtransmatrix)$
    \ENDWHILE
    \RETURN $\newtransmatrix$
  \end{algorithmic}
\end{algorithm}

\begin{algorithm}[t]
  \caption{\label{alg:poweriteration}\yapprox: power iteration approximation of $\vecty_\agroup$}
  \begin{algorithmic}[1]
    \REQUIRE Transition matrix \transmatrix, restart probability $\gamma$, group indicator vector $\indvect{\agroup}$, power iteration exponent~$t$
    \ENSURE Approximation of $\gamma \indvect{\agroup}^T \matrixU^{-1}$
    \STATE $\vecty \gets \indvect{\agroup}$
    \STATE $\matrixU \gets \indvect{\agroup}$ \aris{Is \matrixU a matrix or a vector?}
    \STATE $\indvect{s} \gets \mathbf{1} - \transmatrix \mathbf{1}$
    \FOR {$i = 1$ to $t$}
      \STATE $\matrixU \gets (1 - \gamma) (\transmatrix \matrixU + \indvect{s} \vectv^T \matrixU)$
      \STATE $\vecty \gets \vecty + \matrixU$
    \ENDFOR
    \RETURN $\gamma \vecty$
  \end{algorithmic}
\end{algorithm}

\begin{algorithm}[t]
  \caption{\label{alg:project}\project: projection to feasible set}
  \begin{algorithmic}[1]
    \REQUIRE Transition matrix $\newtransmatrix$
    \ENSURE Projected transition matrix $\newtransmatrix$
    \STATE $\newtransmatrix[i, j] \gets 0$ for $(i, j)$ where $\newtransmatrix[i, j] < 0$
    \STATE $\vects = \newtransmatrix \mathbf{1}$
    \STATE $\newtransmatrix[i, j] \gets \newtransmatrix[i, j] / \vects[i]$ for $(i, j)$ where $\newtransmatrix[i, j] > 0$
  \end{algorithmic}
\end{algorithm}

\spara{Time complexity analysis.} \todo{}

\subsection{Cross-group fair Page\-Rank}

The \fairgd algorithm can be extended to solve the cross-group $\grouppr$-fair Page\-Rank 
problem by adding an outer loop that iterates through all groups. 
The algorithm, dubbed \crossgd, is presented in Algorithm~\ref{alg:crossgd}.

\begin{algorithm}
  
  \caption{\label{alg:crossgd}\crossgd: cross-group fair Page\-Rank projected gradient descent}
  \begin{algorithmic}[1]
    \REQUIRE Transition matrix \transmatrix, restart probability $\gamma$, groups $\vertices_1,\ldots, \vertices_\nogroups$, target Page\-Rank scores $\grouppr_1, \ldots, \grouppr_\nogroups$, learning rate $\alpha$, power iteration exponent~$t$, convergence criterion $\epsilon$
    \ENSURE Optimized transition matrix $\newtransmatrix$
    \STATE $\newtransmatrix \gets \transmatrix$
    \STATE ${\iter} \gets 0$ \aris{\iter is not used}
    \STATE ${\lossprev} \gets \infty$
    \WHILE {$|\losscross(\newtransmatrix, \gamma, \vectv) - {\lossprev}| > \epsilon$}
      \STATE ${\lossprev} \gets \losscross(\newtransmatrix, \gamma, \vectv)$
      \FOR {$\ell = 1$ to $\nogroups$}
        \STATE $\prvector_{\ell} \gets {\pralgo}(\newtransmatrix, \gamma,  \frac{1}{|\vertices_{\ell}|} \indvect{\ell})$
        \FOR {$\agroup = 1$ to $\nogroups$}
          \STATE $\theta_\agroup \gets \indvect{\agroup}^T \prvector_{\ell}$
          \STATE $\vecty_{\agroup} \gets$ ${\yapprox}(\newtransmatrix, \gamma, \indvect{\agroup}, t)$
          \FOR {$(i, j)$ where $\transmatrix[i, j] > 0$}
            \STATE $\newtransmatrix[i, j] \gets \newtransmatrix[i, j] - \alpha \frac{2(1 - \gamma)}{\gamma} (\theta_\agroup - \grouppr_\agroup) \prvector_{\ell}[i] \vecty_{\agroup}[j]$
          \ENDFOR
        \ENDFOR
      \ENDFOR
      \STATE $\newtransmatrix \gets {\project}(\newtransmatrix)$
    \ENDWHILE
    \RETURN $\newtransmatrix$
  \end{algorithmic}
\end{algorithm}

\section{Experimental evaluation  [Haoyun]}

\aris{Algorithms for community detection: (1) Louvain (2) METIS}

% research questions: what are the things that we want to evaluate in our experiments?
% -- comparison with baselines
% -- how the parameters influence the results?
% -- how the methods scale with the size of the data?
% description of the datasets used in the experiments
% methods we try: description of our methods (different variants) and baselines
% what metrics we use for comparison?
% what are the results?
% -- tables and figures

We evaluate the performance of our \fairgd and \crossgd algorithms on real-world networks. In particular, we study the following research questions:


\noindent \textbf{RQ1:} How much can the \fairgd and \crossgd algorithms improve the fairness and cross-group fairness of the PageRank scores? \\
\textbf{RQ2:} When $\grouppr$-fairness and cross-group $\grouppr$-fairness are achieved, how do the \fairgd and \crossgd algorithms compare with the baselines in terms of changes to the transition matrix? \\
\textbf{RQ3:} How do the algorithms scale with the size of the data? \\

The changes to the transition matrix are measured by the relative change in the Frobenius norm of the transition matrix:
\begin{equation}
  \delta = \frac{||\newtransmatrix - \transmatrix||_F}{||\transmatrix||_F}.
\end{equation}

The fairness and cross-group fairness of PageRank scores are measured by the fairness loss and cross-group fairness loss defined in Definitions~\ref{definition:fairness-loss} and \ref{definition:cgfl}.

\subsection{Datasets}

\begin{table}
  \caption{Datasets used in the experiments.}
  \label{table:datasets}
  \begin{tabular}{lrrrc}
    \toprule
    Dataset & \#Vertices & \#Edges & \#Groups & Modularity \\
    \midrule
    Books & 105 & 882 & 3 & 0.4149 \\
    Blogs & 1224 & 19025 & 2 & 0.4111 \\
    Twitter & 18470 & 48365 & 2 & 0.4755 \\
    Slashdot & 82168 & 948464 & 4 & 0.2950 \\
    Google & 875713 & 5105039 & 8 & 0.7420 \\
    \bottomrule
  \end{tabular}
\end{table}

The datasets used in the experiments are as follows:

\noindent \textbf{Books} \footnote{\url{http://www-personal.umich.edu/~mejn/netdata/}}. An undirected network of political books sold by Amazon. Each edge represents frequent co-purchases of two books. The books are divided into three groups: liberal, neutral, and conservative. Since the network is undirected, we convert each edge to two directed edges. \\
\textbf{Blogs} \citep{adamic_political_2005}. A directed network of political blogs. Each edge represents a hyperlink between two blogs. The two vertex groups represent liberal and conservative blogs. \\
\textbf{Twitter} \citep{nr-sigkdd16}. A directed political retweet network. Each edge represents a retweet from one user to another. The two vertex groups are left-leaning and right-leaning users.

The following datasets are not labeled with their group information. We use METIS graph partitioning algorithm \citep{karypis_metis_1997} to partition the vertices into groups. 

\noindent \textbf{Slashdot} \citep{leskovec_community_2009}. A directed network where edges are friend/foe links between users of Slashdot. \\
\noindent \textbf{Google} \citep{leskovec_community_2009}. A directed network where vertices are web pages and edges are hyperlinks between web pages.

\subsection{Baselines}

We compare the \fairgd and \crossgd algorithms with the following baselines:

\noindent \textbf{LFBRN} \citep{tsioutsiouliklis_fairness-aware_2021}. 
A neighborhood-based algorithm that achieves perfect $\phi$-fairness in PageRank scores. \\
\textbf{LFBRU} \citep{tsioutsiouliklis_fairness-aware_2021}. 
A residual-based algorithm that achieves perfect $\phi$-fairness in PageRank scores. \\
\textbf{FairWalk} \citep{rahman_fairwalk_2019}. 
Reweight the edges such that a random walker has equal probability of visiting vertices in each reachable group. \\
\textbf{CrossWalk} \citep{khajehnejad_crosswalk_2022}. 
Based on FairWalk, also take into account the distance of each vertex to other groups in the random walk.

LFBRN and LFBRU are guaranteed to achieve perfect $\phi$-fairness in PageRank scores,
where the restart vector is $\phi$-fair, i.e., $\indvect{g}^T v = \phi_g$.
Other compared methods are assumed to have uniform restart vectors.
FairWalk and CrossWalk are two algorithms that rebalance the random walk 
to achieve fairness in learning graph embeddings, without adding new edges. 
Originally, FairWalk and CrossWalk assign equal weights to each group in the random walk. 
In our experiments, we modify FairWalk and CrossWalk that the probability of 
visiting each group for a random walk is proportional to the target group-wise PageRank scores.
In later experiments, we will show that FairWalk and CrossWalk can improve 
the fairness of PageRank scores but do not guarantee perfect fairness.

\subsection{Results}

\subsubsection{Fair Page\-Rank}

\begin{table*}[htbp]
	\centering
	\caption{Blogs, varying target PageRank ratios. \\
    Initial PageRank ratio: $[0.4835, 0.5165]$.}
	\label{tab:blogs}
	\begin{tabular}{cc|ccccc|ccccc}
		\hline
    \multirow{2}{*}{$\phi[1]$} & \multirow{2}{*}{$\phi[2]$} & \multicolumn{5}{c}{$L$ (Fairness loss)} & \multicolumn{5}{c}{$\delta$ (Relative change)} \\
    \cline{3-7} \cline{8-12}
    & & FairGD & LFPRU & LFPRN & FairWalk & CrossWalk & FairGD & LFPRU & LFPRN & FairWalk & CrossWalk \\

		\hline
    0.1 & 0.9 & 0.03049746 & 0.0000 & 0.0000 & 0.0271 & 0.0752 & 0.2815 & 0.6330 & 0.8178 & 0.5893 & 0.7798 \\
    0.2 & 0.8 & 0.00335592 & 0.0000 & 0.0000 & 0.0165 & 0.0374 & 0.3523 & 0.5628 & 0.7299 & 0.5137 & 0.7239 \\
    0.3 & 0.7 & 0.00000066 & 0.0000 & 0.0000 & 0.0083 & 0.0132 & 0.1016 & 0.5131 & 0.6630 & 0.4544 & 0.6674 \\
    0.4 & 0.6 & 0.00000001 & 0.0000 & 0.0000 & 0.0028 & 0.0012 & 0.0276 & 0.4854 & 0.6239 & 0.4182 & 0.6401 \\
    0.5 & 0.5 & 0.00000057 & 0.0000 & 0.0000 & 0.0001 & 0.0020 & 0.0040 & 0.4811 & 0.6180 & 0.4112 & 0.6596 \\
    0.6 & 0.4 & 0.00000077 & 0.0000 & 0.0000 & 0.0006 & 0.0151 & 0.0347 & 0.5004 & 0.6461 & 0.4350 & 0.6448 \\
    0.7 & 0.3 & 0.00000021 & 0.0000 & 0.0000 & 0.0046 & 0.0444 & 0.2020 & 0.5411 & 0.7042 & 0.4850 & 0.6801 \\
    0.8 & 0.2 & 0.00873229 & 0.0000 & 0.0000 & 0.0122 & 0.0816 & 0.2430 & 0.6005 & 0.7856 & 0.5541 & 0.7661 \\
    0.9 & 0.1 & 0.04409839 & 0.0000 & 0.0000 & 0.0237 & 0.1341 & 0.2575 & 0.6781 & 0.8840 & 0.6362 & 0.8209 \\
    
		\hline
	\end{tabular}
\end{table*}

\begin{table*}[htbp]
	\centering
	\caption{Twitter, varying target PageRank ratios. \\
    Initial PageRank ratio: $[0.4234, 0.5766]$.}
	\label{tab:twitter}
	\begin{tabular}{cc|ccccc|ccccc}
		\hline
    \multirow{2}{*}{$\phi[1]$} & \multirow{2}{*}{$\phi[2]$} & \multicolumn{5}{c}{$L$ (Fairness loss)} & \multicolumn{5}{c}{$\delta$ (Relative change)} \\
    \cline{3-7} \cline{8-12}
    & & FairGD & LFPRU & LFPRN & FairWalk & CrossWalk & FairGD & LFPRU & LFPRN & FairWalk & CrossWalk \\

		\hline
    0.1 &	0.9 & 0.12347248 & 0.0 & 0.0 & 0.0860 & 0.0987 & 0.5182 & 0.6097 & 0.6407 & 0.2312 & 0.5413 \\
    0.2 &	0.8 & 0.04543222 & 0.0 & 0.0 & 0.0395 & 0.0470 & 0.5011 & 0.5548 & 0.5822 & 0.1986 & 0.5223 \\
    0.3 &	0.7 & 0.00712386 & 0.0 & 0.0 & 0.0109 & 0.0142 & 0.4267 & 0.5157 & 0.5388 & 0.1699 & 0.4932 \\
    0.4 &	0.6 & 0.00000043 & 0.0 & 0.0 & 0.0001 & 0.0004 & 0.0571 & 0.4955 & 0.5145 & 0.1475 & 0.5135 \\
    0.5 &	0.5 & 0.00000083 & 0.0 & 0.0 & 0.0072 & 0.0060 & 0.3935 & 0.4958 & 0.5119 & 0.1345 & 0.5069 \\
    0.6 &	0.4 & 0.00628326 & 0.0 & 0.0 & 0.0321 & 0.0312 & 0.7963 & 0.5168 & 0.5314 & 0.1338 & 0.4964 \\
    0.7 &	0.3 & 0.04881410 & 0.0 & 0.0 & 0.0749 & 0.0746 & 0.7967 & 0.5560 & 0.5707 & 0.1455 & 0.5039 \\
    0.8 &	0.2 & 0.13094114 & 0.0 & 0.0 & 0.1356 & 0.1383 & 0.7987 & 0.6101 & 0.6260 & 0.1670 & 0.5274 \\
    0.9 &	0.1 & 0.25326871 & 0.0 & 0.0 & 0.2143 & 0.2209 & 0.7987 & 0.6761 & 0.6937 & 0.1951 & 0.5188 \\

		\hline
	\end{tabular}
\end{table*}

First, we study the condition under which $\phi$-fairness and cross-group $\phi$-fairness can be achieved by the FairGD and CrossGD algorithms. We run the FairGD on the Blogs and Twitter with a uniform restart vector. We vary the target PageRank scores of the two groups and record the fairness loss $L$ and the relative change to the transition matrix $\delta$ with FairGD and baselines. The results are shown in Table~\ref{tab:blogs} and Table~\ref{tab:twitter}.

\haoyun{We can add a null baseline that does not change the transition matrix, to show the initial fairness loss.}

\haoyun{The losses are often very close to zero. We might use a logarithmic scale.}

\haoyun{The results for varying target PageRank ratios can be shown in figures. 
Each figure shows a metric (fairness loss or relative change) with respect to the target PageRank ratio.}

% For Blogs dataset, we find that the fairness loss with FairGD algorithm 
% is approximately zero when the target PageRank ratios range from 0.3 to 0.7 
% for the first group, indicating that the $\phi$-fairness in PageRank scores is achieved. 
% In these cases, FairGD outperforms LFBRU and LFBRN in terms of 
% the relative change to the transition matrix. 
% Especially when the target PageRank ratios are close to 0.5, 
% the relative change to the transition matrix with FairGD 
% is significantly smaller than that with LFBRU and LFBRN.

% For Twitter dataset, the FairGD algorithm achieves $\phi$-fairness in PageRank scores 
% when the target PageRank ratios range from 0.4 to 0.5 for the first group. 
% Again, FairGD outperforms LFBRU and LFBRN in terms of 
% the relative change to the transition matrix.

The Blogs and Twitter datasets contain two groups.
We vary the target Page\-Rank scores of the two groups
from 0.2 to 0.8 and from 0.8 to 0.2, respectively.
As shown in Table~\ref{tab:blogs} and Table~\ref{tab:twitter},
when the target Page\-Rank scores are close to the initial Page\-Rank scores,
FairGD has a significantly smaller relative change to the transition matrix
than LFBRU and LFBRN, while having an almost zero fairness loss.
In these cases, FairGD also outperforms FairWalk and CrossWalk in terms of the Fairness loss.

The other datasets contain more than two groups. 
We set the target PageRank scores $\phi$ proportional to the number of vertices in each group, 
indicating the demographic parity. 
From Table~\ref{tab:multi}, we observe that the FairGD algorithm achieves
$\phi$-fairness in PageRank scores for all datasets,
with the fairness loss close to zero.
In terms of the relative change to the transition matrix,
FairGD outperforms all baseline methods in all datasets.
In terms of the fairness loss, FairGD also outperforms FairWalk and CrossWalk.

% We notice the relationship between the performance of FairGD and the modularity of the network.

\haoyun{Is there a better way to show the tables?}

\begin{table*}[htbp]
	\centering
	\caption{Multi-group datasets.}
	\label{tab:multi}
	\begin{tabular}{l|ccccc|ccccc}
		\hline
    \multirow{2}{*}{Dataset} & \multicolumn{5}{c}{$L$ (Fairness loss)} & \multicolumn{5}{c}{$\delta$ (Relative change)} \\
    \cline{2-6} \cline{7-11}
    & FairGD & LFPRU & LFPRN & FairWalk & CrossWalk & FairGD & LFPRU & LFPRN & FairWalk & CrossWalk \\

		\hline

    Books & 0.00000099 & 0.0 & 0.0 & 0.0003 & 0.0005 & 0.1406 & 0.6053 & 0.6101 & 0.1803 & 0.6510 \\

    Slashdot & 0.00000075 & 0.0 & 0.0 & 0.0003 & 0.0010 & 0.0601 & 0.6921 & 0.6929 & 0.2775 & 0.3287 \\

    Google & 0.00000081 & 0.0 & 0.0 & 0.000007 & 0.000006 & 0.0043 & 0.8623 & 0.8560 & 0.3095 & 0.4835 \\

		\hline
	\end{tabular}
\end{table*}

\haoyun{For google dataset, the initial PageRank scores are almost fair. Tried different numbers of partition but still fair.}

\subsubsection{Cross-group Fair Page\-Rank}

We now study \crossgd algorithm.

\begin{table*}[htbp]
	\centering
	\caption{Cross-group fairness.}
	\label{tab:cg}
	\begin{tabular}{l|ccccc|ccccc}
		\hline
    \multirow{2}{*}{Dataset} & \multicolumn{5}{c}{$L$ (Fairness loss)} & \multicolumn{5}{c}{$\delta$ (Relative change)} \\
    \cline{2-6} \cline{7-11}
    & FairGD & LFPRU & LFPRN & FairWalk & CrossWalk & FairGD & LFPRU & LFPRN & FairWalk & CrossWalk \\

		\hline

    Books & 0.1074 & 0.0000 & 0.0000 & 0.1093 & 0.1208 & 0.6118 & NA & NA & 0.1803 & 0.6510 \\

    Blogs & 0.0236 & 0.0000 & 0.0000 & 0.0431 & 0.1277 & 0.8882 & NA & NA & 0.4101 & 0.6316 \\

    % Twitter & 

    % Slashdot & 0.00000075 & 0.0 & 0.0 & 0.0003 & 0.0010 & 0.0601 & 0.6921 & 0.6929 & 0.2775 & 0.3287 \\

    % Google & 0.00000081 & 0.0 & 0.0 & 0.000007 & 0.000006 & 0.0043 & 0.8623 & 0.8560 & 0.3095 & 0.4835 \\

		\hline
	\end{tabular}
\end{table*}

\begin{table*}[htbp]
  \centering
	\caption{Fairness loss.}
  \begin{tabular}{l|ccccc}
    \hline
    & FairGD & LFPRU & LFPRN & FairWalk & CrossWalk \\
    \hline
    Books & 0.00000099 & 0.0 & 0.0 & 0.00026933 & 0.00051765 \\
    % Blogs & 0.00000075 & 0.0 & 0.0 & 0.0003 & 0.0010 \\
    % Twitter & 0.00000081 & 0.0 & 0.0 & 0.000007 & 0.000006 \\
    Slashdot & 0.00000075 & 0.0 & 0.0 & 0.00027270 & 0.00100066 \\
    Google & 0.00000081 & 0.0 & 0.0 & 0.00000701 & 0.00000558 \\
    \hline
    
  \end{tabular}
\end{table*}

\begin{table*}[htbp]
  \centering
	\caption{Fairness loss \loss.}
  \begin{tabular}{l|ccccc}
    \hline
    & FairGD & LFPRU & LFPRN & FairWalk & CrossWalk \\
    \hline
    Books & 0.00000099 & 0.0 & 0.0 & 0.00026933 & 0.00051765 \\
    Blogs & 0.00000075 & 0.0 & 0.0 & 0.00047937 & 0.00073687 \\
    Twitter & 0.00000285 & 0.0 & 0.0 & 0.00056446 & 0.00105136 \\
    Slashdot & 0.00000075 & 0.0 & 0.0 & 0.00027270 & 0.00100066 \\
    Google & 0.00000081 & 0.0 & 0.0 & 0.00000701 & 0.00000558 \\
    \hline
    
  \end{tabular}
\end{table*}

% \begin{table*}[htbp]
%   \centering
% 	\caption{Change to transition matrix.}
%   \begin{tabular}{l|ccccc}
%     \hline
%     & FairGD & LFPRU & LFPRN & FairWalk & CrossWalk \\
%     \hline
%     Books & 0.00000099 & 0.0 & 0.0 & 0.00026933 & 0.00051765 \\
%     Blogs & 0.00000075 & 0.0 & 0.0 & 0.00047937 & 0.00073687 \\
%     Twitter & 0.00000285 & 0.0 & 0.0 & 0.00056446 & 0.00105136 \\
%     Slashdot & 0.00000075 & 0.0 & 0.0 & 0.00027270 & 0.00100066 \\
%     Google & 0.00000081 & 0.0 & 0.0 & 0.00000701 & 0.00000558 \\
%     \hline
    
%   \end{tabular}
% \end{table*}

%%
%% The acknowledgments section is defined using the "acks" environment
%% (and NOT an unnumbered section). This ensures the proper
%% identification of the section in the article metadata, and the
%% consistent spelling of the heading.
\begin{acks}
\end{acks}

%%
%% The next two lines define the bibliography style to be used, and
%% the bibliography file.
\bibliographystyle{ACM-Reference-Format}

\bibliography{refs.bib}


%%
%% If your work has an appendix, this is the place to put it.
\appendix




\end{document}
\endinput
%%
%% End of file `sample-sigconf-authordraft.tex'.
